\documentclass[12pt,oneside]{book}

% =====================================================
% METADATOS
% =====================================================
\newcommand{\universidad}{Universidad de Buenos Aires (UBA)}
\newcommand{\facultad}{Facultad de Derecho}
\newcommand{\carrera}{Carrera de Abogacía}
\newcommand{\materia}{Práctica Profesional -- Derecho Procesal Civil y Comercial}
\newcommand{\titulo}{Prueba Testimonial, Técnica de Audiencia y Nulidades Procesales}
\newcommand{\autor}{Isaac Edgar Camacho Ocampo}
\newcommand{\anio}{2026}

% =====================================================
% PAQUETES
% =====================================================
\usepackage[utf8]{inputenc}
\usepackage[T1]{fontenc}
\usepackage[spanish,es-nodecimaldot]{babel}

\usepackage[
    top=2.5cm,
    bottom=2.5cm,
    left=3cm,
    right=2.5cm,
    headheight=15pt
]{geometry}

\usepackage{amsmath}
\usepackage{amssymb}
\usepackage{mathtools}

\usepackage{graphicx}
\usepackage{float}
\usepackage{caption}
\usepackage{subcaption}

\usepackage{booktabs}
\usepackage{array}
\usepackage{longtable}
\usepackage{tabularx}

\usepackage{enumitem}

\usepackage{xcolor}
\usepackage[most]{tcolorbox}

\usepackage{fancyhdr}
\usepackage{ragged2e}

\graphicspath{
    {./img/}
    {./graficos/}
    {./figuras/}
}

\usepackage[
    colorlinks=true,
    linkcolor=blue!50!black,
    citecolor=green!40!black,
    urlcolor=blue!60!black,
    pdftitle={\titulo},
    pdfauthor={\autor},
    pdfsubject={Apuntes universitarios},
    bookmarks=true
]{hyperref}

% =====================================================
% CONFIGURACIÓN GENERAL
% =====================================================
\setcounter{tocdepth}{2}

\setlength{\parindent}{0pt}
\setlength{\parskip}{0.5em}

\setlist[itemize]{
    topsep=0.3em,
    itemsep=0.2em
}
\setlist[enumerate]{
    topsep=0.3em,
    itemsep=0.2em
}

\clubpenalty=10000
\widowpenalty=10000

% =====================================================
% COMANDOS MATEMÁTICOS
% =====================================================
\newcommand{\abs}[1]{\left|#1\right|}
\newcommand{\norm}[1]{\left\lVert#1\right\rVert}
\newcommand{\vect}[1]{\mathbf{#1}}
\newcommand{\deriv}[2]{\frac{d#1}{d#2}}
\newcommand{\pderiv}[2]{\frac{\partial#1}{\partial#2}}

% =====================================================
% ENCABEZADOS Y PIES
% =====================================================
\pagestyle{fancy}
\fancyhf{}
\fancyhead[L]{\nouppercase{\leftmark}}
\fancyhead[R]{\materia}
\fancyfoot[C]{\thepage}
\renewcommand{\headrulewidth}{0.4pt}
\renewcommand{\footrulewidth}{0pt}

\fancypagestyle{plain}{
    \fancyhf{}
    \fancyfoot[C]{\thepage}
    \renewcommand{\headrulewidth}{0pt}
}

% =====================================================
% CAJAS ACADÉMICAS
% =====================================================
\newtcolorbox{definition}[1]{
    enhanced,
    breakable,
    title={Definición: #1},
    fonttitle=\bfseries,
    colback=blue!3,
    colframe=blue!50!black,
    boxrule=0.8pt,
    arc=2mm
}

\newtcolorbox{important}{
    enhanced,
    breakable,
    title={Importante},
    fonttitle=\bfseries,
    colback=yellow!8,
    colframe=orange!70!black,
    boxrule=0.8pt,
    arc=2mm
}

\newtcolorbox{example}[1]{
    enhanced,
    breakable,
    title={Ejemplo: #1},
    fonttitle=\bfseries,
    colback=green!4,
    colframe=green!40!black,
    boxrule=0.8pt,
    arc=2mm
}

\newtcolorbox{formula}[1]{
    enhanced,
    breakable,
    title={Fórmula / Regla: #1},
    fonttitle=\bfseries,
    colback=purple!4,
    colframe=purple!50!black,
    boxrule=0.8pt,
    arc=2mm
}

\newtcolorbox{exercise}[1]{
    enhanced,
    breakable,
    title={Ejercicio: #1},
    fonttitle=\bfseries,
    colback=gray!5,
    colframe=gray!60!black,
    boxrule=0.8pt,
    arc=2mm
}

\newtcolorbox{note}{
    enhanced,
    breakable,
    title={Nota},
    fonttitle=\bfseries,
    colback=white,
    colframe=gray!50,
    boxrule=0.6pt,
    arc=2mm
}

% =====================================================
% DOCUMENTO
% =====================================================
\begin{document}

% -----------------------------------------------------
% PORTADA
% -----------------------------------------------------
\begin{titlepage}
\thispagestyle{empty}

\begin{center}

\vspace*{1cm}

{\large \textsc{\universidad}}

\vspace{0.2cm}

{\large \textsc{\facultad}}

\vspace{0.2cm}

{\large \carrera}

\vspace{3cm}

{\Huge \textbf{\titulo}}

\vspace{1cm}

{\LARGE \materia}

\vspace{0.8cm}

\textit{Comprender el porqué de cada norma procesal es lo que distingue a quien la aplica de quien apenas la recita.}

\vspace{1.2cm}

\rule{0.70\textwidth}{0.5pt}

\vspace{0.5cm}

{\large Apuntes de estudio}

\vspace{0.5cm}

\rule{0.70\textwidth}{0.5pt}

\vfill

{\large \autor}

\vspace{0.3cm}

{\large \anio}

\end{center}

\vfill

\begin{flushleft}
\textit{Este es un modesto aporte a los estudiantes de ``Práctica Profesional -- Derecho Procesal Civil y Comercial'' no pretende sustituir el material oficial de la materia.}
\end{flushleft}

\end{titlepage}

% -----------------------------------------------------
% ÍNDICE
% -----------------------------------------------------
\tableofcontents

% =====================================================
% PRESENTACIÓN DE LA CLASE
% =====================================================
\chapter*{Presentación de la clase}
\addcontentsline{toc}{chapter}{Presentación de la clase}

La clase corresponde a la etapa de práctica profesional de Derecho Procesal Civil y Comercial. Se desarrolla como una devolución sobre una audiencia testimonial simulada por los propios estudiantes, seguida de una explicación teórica sobre la prueba testimonial, la técnica de interrogatorio y los recursos disponibles en la audiencia. La clase cierra con un caso real relatado por el docente sobre un embargo trabado por una notificación defectuosa, y con pautas para la evaluación final.

\section*{Temas tratados en esta clase}
\addcontentsline{toc}{section}{Temas tratados en esta clase}

\begin{itemize}
    \item Revisión crítica de una audiencia testimonial simulada y de la dinámica entre las partes, sus letrados y el/la juzgador(a).
    \item Teoría de la prueba testimonial: el concepto de testigo, sus requisitos y las cargas procesales de la parte que lo ofrece.
    \item Técnica de interrogatorio: preguntas abiertas frente a preguntas indicativas, objeciones y su resolución en la audiencia.
    \item Los recursos frente a las resoluciones sobre prueba: irrecurribilidad, reposición y nulidad procesal.
    \item Un caso real de embargo originado en una notificación defectuosa: tipos de domicilio y estrategia de nulidad.
    \item Pautas para la evaluación final y una reflexión sobre las habilidades profesionales que el ejercicio del derecho exige hoy.
\end{itemize}

\section*{Utilidad de estos contenidos}
\addcontentsline{toc}{section}{Utilidad de estos contenidos}

Esta clase reconstruye, paso a paso, lo que ocurre en una audiencia de prueba testimonial y en la resolución de un embargo mal trabado: notificar a tiempo, formular la pregunta correcta, objetar cuando corresponde y detectar una notificación nula antes de que el cliente vea afectada su cuenta bancaria.

Más allá de la aplicación estrictamente procesal, los contenidos entrenan una habilidad transferible: la de hacer la pregunta correcta, verificar los datos antes de actuar y no dar por cierto lo que solo parece serlo. La diferencia entre un litigante eficaz y uno que expone a su cliente a un daño evitable no suele estar en el conocimiento de la norma, sino en el hábito de aplicarla con método y en el momento oportuno.

% =====================================================
% BLOQUE 1
% =====================================================
\chapter{De la audiencia a la estrategia}

\section{La audiencia en revisión}

La clase comienza con la revisión de la audiencia testimonial simulada que los estudiantes habían llevado adelante en un encuentro anterior, reconstruyendo qué hechos habían quedado controvertidos y qué había ocurrido durante la producción de la prueba.

Del análisis surge que la parte actora no supo explicar con precisión para qué había citado a sus propios testigos, es decir, no tenía claro qué hecho concreto debía acreditar cada uno. Se distingue con claridad la responsabilidad involucrada: ese defecto no constituye un error del testigo, sino una falencia de preparación del letrado, quien debió reunirse previamente con él y explicarle con precisión el objeto de su declaración. Se señala, además, que la falta de coordinación previa entre el letrado y su cliente ---ni siquiera estando presentes en el mismo espacio antes de la audiencia--- impidió cualquier ajuste de último momento.

\begin{important}
La preparación previa del testigo ---explicarle con precisión sobre qué hechos va a declarar--- no es una formalidad: es lo que evita que la parte y su abogado queden expuestos frente al juzgador en plena audiencia. Como se remarca en la clase: \textit{``prepararlo en buen sentido es decirle: nosotros lo llamamos para que acredite tal cosa, simplemente eso''}.
\end{important}

\section{El juez, el abogado y la parte}

El segundo punto de la revisión gira en torno al modo en que la jueza condujo la audiencia. Se advierte que, en un momento, la magistrada se dirigió directamente a la parte actora para interrogarla, en lugar de canalizar la pregunta a través de su letrado, lo que generó cierta incomodidad por tratarse de un rol que debe ser imparcial.

Sobre este punto se aclara que interrogar directamente a la parte no constituye una irregularidad, sino una facultad del tribunal: \textit{``a la parte no solo no está mal interrogarla directamente, sino que no está impedido; es una cuestión de estilo de esta jueza. Está dentro de las posibilidades, con lo cual ustedes tienen que saber que su cliente, por mínimo que sea, tiene que estar preparado porque esto puede suceder''}.

Respecto de la familiaridad percibida entre la jueza ---o su secretaria--- y la parte actora, se relativiza la objeción: al no haber estado presente la otra parte en ese cruce, no llegó a generarse una situación de desigualdad procesal concreta, aunque como principio ese trato debería evitarse para resguardar la imparcialidad.

\section{Caducidad de la prueba testimonial por incomparecencia}

Se plantea un segundo caso, ocurrido en una audiencia ante la misma jueza en un día anterior: al momento de la audiencia testimonial, la parte que había ofrecido a los testigos no acreditó haberlos notificado.

La conclusión a la que se arriba es la siguiente: si los testigos no fueron notificados ---en el ejemplo, mediante carta documento--- y por eso no comparecen, la parte contraria debe pedir en el acto la caducidad de la prueba testimonial ofrecida, invocando el artículo pertinente del Código Procesal.

\begin{important}
\textbf{Norma citada en clase --- artículo 432 del Código Procesal.} Fórmula empleada en la audiencia simulada: \textit{``solicito la caducidad de la prueba testimonial ofrecida por la parte demandada, por el artículo 432, por no haber dado cumplimiento a la notificación; y si quieren quedar bien, hasta pueden citar el artículo en el acto''}.

En ese supuesto no corresponde ``impugnar'' nada ---porque no hay ningún acto procesal que impugnar--- sino solicitar directamente la caducidad de la prueba por incomparecencia imputable a la parte que la ofreció y no cumplió con la carga de notificar.
\end{important}

% =====================================================
% BLOQUE 2
% =====================================================
\chapter{La prueba testimonial: teoría y cargas}

\section{El testigo: concepto y requisitos}

A partir del episodio anterior se construye la teoría general de la prueba testimonial.

\begin{definition}{Testigo}
Es quien conoce por sus propios sentidos ---porque lo vio, lo escuchó o lo percibió de algún modo--- un hecho relevante para el proceso, y lo relata de manera directa, sin que se trate de algo que le hayan contado a él (lo que lo distinguiría de un testigo de oídas o ``de referencia'').
\end{definition}

Sobre los requisitos, se destaca especialmente la \textbf{imparcialidad}: el testigo no debe tener un interés personal en el resultado del pleito. Cualquiera de las partes puede ofrecer testigos, en tanto reúnan esas condiciones.

\section{Cargas procesales al ofrecer un testigo}

Se enumeran las cargas que pesan sobre la parte que ofrece un testigo:

\begin{table}[H]
\centering
\begin{tabularx}{\textwidth}{
    >{\raggedright\arraybackslash}p{0.32\textwidth}
    >{\raggedright\arraybackslash}X
}
\toprule
\textbf{Momento procesal} & \textbf{Carga de la parte que ofrece el testigo} \\
\midrule
Ofrecimiento del testigo &
Identificarlo (nombre y, en la práctica, DNI, aunque el Código no lo exige literalmente así) e indicar sobre qué hechos declarará. \\
Después de ofrecido &
Asegurar su comparecencia: en principio el testigo es citado por el juzgado, salvo que la parte que lo propuso asuma la carga de hacerlo comparecer a la audiencia. \\
Al redactar la demanda &
Precisar respecto de qué hechos concretos se cita a declarar a cada testigo (no es obligatorio incluir ya el cuestionario, pero sí el objeto de la declaración). \\
\bottomrule
\end{tabularx}
\end{table}

\begin{important}
\textbf{Norma citada en clase.} Regla general del Código sobre la citación del testigo: ``el testigo será citado por el juzgado, salvo cuando la parte que lo propuso asuma la carga de hacerlo comparecer a la audiencia''.
\end{important}

Se subraya por qué es necesario indicar el objeto de cada testimonio: aunque los jueces, por el principio de amplitud probatoria, rara vez desestiman un testigo por esta omisión, la contraparte tiene derecho a saber con anticipación a qué viene cada testigo, en resguardo del derecho de defensa y del principio de bilateralidad.

\section{Preparación del testigo y construcción del cuestionario}

Antes de ofrecer un testigo es necesario tomar contacto con él y confirmar que efectivamente tiene conocimiento directo de los hechos; suele ser el propio cliente quien identifica a esas personas. A partir de allí, la estrategia probatoria se organiza en etapas: primero se definen qué hechos hay que probar; luego, a través de qué medios ---en este caso, testigos---; y finalmente se redacta el cuestionario que se le formulará a cada uno.

\begin{important}
Regla práctica no escrita en el Código: ``prepárense ustedes respecto de sus testigos, no se encuentren con una sorpresa el día en que van a declarar y no sepan de qué se trata''.
\end{important}

% =====================================================
% BLOQUE 3
% =====================================================
\chapter{Técnica de interrogatorio en audiencia}

\section{Preguntas abiertas frente a preguntas indicativas}

Las preguntas al testigo deben ser \textbf{preguntas abiertas}, que no puedan contestarse con un simple ``sí'' o ``no''. Las preguntas cerradas tienden a sugerir la respuesta ---son \textbf{preguntas indicativas}--- y deben ser objetadas.

\begin{example}{Simulación de una pregunta indicativa}
Se dramatiza la interrogación de un testigo con la pregunta: ``¿el semáforo estaba en rojo cuando se produjo el choque?''. Esta pregunta debe ser objetada por indicativa, porque no pide que el testigo narre lo que percibió, sino que confirme una versión ya armada dentro de la propia pregunta.

La forma correcta de objetar es solicitar la palabra y decir, por ejemplo: ``me opongo, la pregunta es indicativa'', pidiendo su reformulación.

La pregunta correcta, en cambio, debe ser abierta: ``¿qué sucedió?'' o ``¿qué vio?'', dejando que sea el testigo quien narre los hechos con sus propias palabras.
\end{example}

Esta técnica no es exclusiva del proceso civil: en materia penal el interrogatorio suele ser todavía más estrictamente abierto (``¿qué hizo?, ¿dónde estaba?, ¿tenía el arma?''), evitando cualquier formulación que insinúe la respuesta.

\section{Resolución de objeciones y reformulación}

Ante una discusión entre las partes sobre si una pregunta es o no indicativa, la dinámica habitual de la audiencia es la siguiente: primero se retira al testigo de la sala; luego quien preside el acto invita a las partes a reformular la pregunta de común acuerdo. Si no llegan a un acuerdo, el tribunal tiene dos caminos: reformular la pregunta por su propia decisión, o directamente disponer que esa pregunta no se formule.

\section{Recursos contra las resoluciones sobre prueba}

Se plantea qué remedio procesal tiene la parte disconforme con una decisión sobre la prueba.

\begin{important}
\textbf{Norma citada en clase.} ``Todas las cuestiones relativas a la producción y denegación de la prueba son irrecurribles, son inapelables''; la vía que sí queda disponible es la reposición: ``por efecto diferido, no; reposición sí''.
\end{important}

El recurso de reposición debe interponerse en el mismo acto ---no después--- y con fundamento inmediato (por ejemplo: ``pido reposición porque esta pregunta es determinante y el criterio para dejarla sin efecto es arbitrario''). El tribunal, antes de resolver, corre traslado a la contraparte, que puede sostener su posición u oponerse a la reposición. Recién con ambas posiciones escuchadas, se resuelve ``hágase lugar'' o ``no hágase lugar'', con un fundamento breve. Si el tribunal no llegara a escuchar la posición de alguna de las partes antes de resolver, la consecuencia posible sería plantear la nulidad de lo actuado.

\begin{important}
La reposición se pide en el momento, se funda en el momento y se resuelve después de escuchar a ambas partes. No hacerlo en el acto significa perder la posibilidad de revertir la decisión sobre la prueba, dado que estas resoluciones son inapelables.
\end{important}

% =====================================================
% BLOQUE 4
% =====================================================
\chapter{Domicilios y nulidad procesal: un caso real}

\section{Los tipos de domicilio en el proceso}

A partir de un caso real ---un embargo originado en una cuenta bancaria--- se revisan las distintas clases de domicilio relevantes en el proceso.

\begin{table}[H]
\centering
\begin{tabularx}{\textwidth}{
    >{\raggedright\arraybackslash}p{0.24\textwidth}
    >{\raggedright\arraybackslash}X
}
\toprule
\textbf{Tipo de domicilio} & \textbf{Descripción} \\
\midrule
Real & Aquel en el que la persona efectivamente vive, o donde la sociedad está legalmente constituida. \\
Denunciado & El que la parte actora indica como domicilio real de la contraria al iniciar la demanda, sin que todavía esté verificado. \\
Constituido & El que cada parte fija a los fines procesales dentro de un expediente determinado; debe estar dentro del radio del juzgado y hoy se le suma el domicilio electrónico. \\
Electrónico & El domicilio procesal constituido en el sistema informático del Poder Judicial, que se agrega al domicilio constituido tradicional. \\
Fiscal & El que la persona o la sociedad tiene registrado ante el organismo de recaudación (AFIP/ARCA) o ante la provincia; no debe confundirse con el domicilio real ni tiene, en principio, efectos procesales sobre la notificación. \\
Legal & El que la ley asigna a ciertas personas por su función; como ejemplo excepcional se menciona el del nuncio apostólico. \\
\bottomrule
\end{tabularx}
\end{table}

\begin{important}
En el caso de una sociedad demandada, lo que debe coincidir con el domicilio al que se notifica es el domicilio real registrado ---por ejemplo, el que figura ante la Inspección General de Justicia (IGJ) en el contrato social--- y no cualquier otro domicilio que pueda figurar en otros registros.
\end{important}

\section{El caso del embargo por notificación defectuosa}

\begin{example}{Caso real relatado en clase}
El cliente es una sociedad anónima (S.A.) demandada en un juicio laboral. La cédula de notificación fue dirigida a una calle Rivadavia, domicilio que ---según el contrato social presentado ante la IGJ--- no coincide con el domicilio real de la sociedad.

La cédula fue recibida por una persona que se identificó como empleada y firmó el acuse de recibo, pero la sociedad demandada no tiene ninguna relación con ese domicilio: probablemente se trate del domicilio particular del dueño de la empresa, y la persona que firmó tenga con él una relación de índole personal, ajena por completo a la S.A.

Como consecuencia de esa notificación mal dirigida, el proceso siguió su curso sin que la demandada se enterara: hubo rebeldía y, finalmente, un embargo de aproximadamente cien millones de pesos sobre la cuenta del cliente.
\end{example}

La pregunta central a resolver es qué puede hacer el abogado de la demandada frente a este escenario. Se descarta primero la apelación, porque los plazos para apelar ya están vencidos: el embargo ya se trabó, de modo que esa vía ya no está disponible. El camino identificado es el planteo de nulidad procesal, con precisión técnica sobre qué acto exactamente debe recaer esa nulidad.

\section{La nulidad en cadena: notificación, rebeldía y embargo}

\begin{important}
\textbf{Sobre qué se pide la nulidad.} No corresponde pedir la ``nulidad de la sentencia'' en abstracto, sino, en primer lugar, la nulidad de la notificación defectuosa ---el acto que originó todo el resto del vicio---.

Como consecuencia de esa nulidad, caen también los actos posteriores que dependían de ella: la declaración de rebeldía y, en cadena, la sentencia dictada en rebeldía.

Respecto del embargo, la precisión es todavía más fina: el embargo en sí mismo no es nulo, porque fue trabado conforme a derecho a partir de la sentencia vigente en ese momento. Lo que corresponde pedir, una vez declarada la nulidad de lo anterior, es el levantamiento del embargo, como consecuencia de que ya no subsiste la sentencia que le dio origen.
\end{important}

La vía procesal concreta para plantear todo esto es un escrito con un sumario o título que identifique con precisión lo que se pide (nulidad de la notificación y, como consecuencia, de los actos posteriores, con solicitud de levantamiento de embargo), acompañado de la prueba que acredite el vicio ---en este caso, el contrato social inscripto ante la IGJ, que demuestra cuál es el verdadero domicilio real de la sociedad---.

Como segundo ejemplo del mismo problema, se relata el caso de una clienta de un pueblo de Córdoba, clienta de un mismo banco desde hacía casi cincuenta años, a quien un día el banco le informó que su cuenta estaba embargada. El origen era un juicio ejecutivo en el que también se la había notificado de manera ficta en un domicilio en el que ella no vivía, y la solución que debió acreditarse fue sustancialmente la misma: la nulidad de esa notificación defectuosa.

\begin{important}
Verificar el domicilio al que se dirigió cada notificación ---antes de cualquier otro análisis--- es el primer paso frente a cualquier acto procesal sorpresivo (un embargo, una rebeldía, una sentencia desconocida). Muchas veces el vicio no está en el fondo del reclamo, sino en la forma en que se llevó adelante la notificación.
\end{important}

% =====================================================
% BLOQUE 5
% =====================================================
\chapter{Cierre: evaluación y formación profesional}

\section{Pautas para la evaluación final}

\begin{note}
\begin{itemize}
    \item Se permite tener el Código en papel; no se permite ningún dispositivo electrónico. Quien no lleve el Código impreso rendirá sin esa herramienta de consulta.
    \item La letra debe ser legible, sin necesidad de recurrir a un perito calígrafo para descifrarla.
    \item Las preguntas van a ser simples y las respuestas deben ser precisas: ``si les preguntan, traten de contestar; si después quieren agregar algo más, piénsenlo''.
    \item Se recomienda leer todas las preguntas primero para administrar el tiempo, y responder en el orden que a cada estudiante le resulte más cómodo, sin necesidad de seguir el orden del cuestionario.
    \item Quienes deban rendir en instancia de recuperatorio deberán responder, además, un contenido adicional.
\end{itemize}
\end{note}

\section{Habilidades para el abogado del futuro}

Como cierre de la clase se propone una reflexión de tipo profesional más amplia, vinculada con el sentido de todo lo trabajado durante el curso, a partir de la pregunta: ``¿qué es lo que nos puede servir para hoy y, probablemente, dentro de diez años? ¿cuáles son las habilidades que tenemos que explotar?''.

La respuesta parte de una constatación: buena parte de las formaciones que eran útiles hace diez o veinte años hoy quedaron obsoletas, y el conocimiento memorizado es, en ese sentido, relativo: ``apretamos un botón y la información está''. Sin embargo, lo que no se resuelve apretando un botón es otra cosa: ``cómo manejar esa información todavía es algo que lo hacemos los seres humanos, porque estamos corriendo gravísimos riesgos con eso de apretar el botón y creer que con eso estamos resolviendo las cosas''.

Ese manejo criterioso de la información ---no la acumulación de datos--- es, en definitiva, el verdadero entrenamiento que la materia buscó dar a lo largo del curso.

% =====================================================
% CORNELL
% =====================================================
\chapter{Cornell: síntesis de la clase}

La siguiente tabla organiza los conceptos centrales de la clase en formato de notas Cornell: preguntas o conceptos clave a la izquierda y su respuesta o desarrollo a la derecha, seguidas de una síntesis final integradora.

\begin{longtable}{
    >{\raggedright\arraybackslash}p{0.30\textwidth}
    >{\raggedright\arraybackslash}p{0.60\textwidth}
}
\toprule
\textbf{Preguntas / palabras clave} & \textbf{Notas / respuestas} \\
\midrule
\endfirsthead
\toprule
\textbf{Preguntas / palabras clave} & \textbf{Notas / respuestas} \\
\midrule
\endhead

\textbf{¿Qué carga tiene la parte que ofrece un testigo si este no comparece por falta de notificación?} &
Ninguna consecuencia le corresponde a esa parte si notificó correctamente. Pero si el testigo no fue notificado por su culpa, la contraria puede pedir en el acto, con base en el artículo 432, la caducidad de la prueba testimonial ofrecida. \\
\addlinespace[0.5em]

\textbf{¿Puede el juez interrogar directamente a la parte, y no solo a través de su abogado?} &
Sí. No está prohibido ni impedido; es una facultad del tribunal y una cuestión de estilo de cada juez o jueza. Por eso el abogado debe preparar a su cliente para esa posibilidad. \\
\addlinespace[0.5em]

\textbf{¿Qué es un testigo y qué requisito es central para que su declaración tenga valor?} &
Es quien conoce un hecho por percepción directa de sus propios sentidos y lo relata de manera directa. El requisito central es la imparcialidad: no debe tener un interés personal en el resultado del proceso. \\
\addlinespace[0.5em]

\textbf{¿Qué debe indicarse al ofrecer un testigo, además de su identidad?} &
Debe precisarse sobre qué hechos concretos va a declarar, en resguardo del derecho de defensa y de la bilateralidad, aunque los jueces rara vez desestiman el testimonio por esta sola omisión. \\
\addlinespace[0.5em]

\textbf{¿Cuál es la diferencia entre una pregunta abierta y una pregunta indicativa?} &
La abierta (``¿qué sucedió?'', ``¿qué vio?'') deja que el testigo narre libremente. La indicativa (``¿el semáforo estaba en rojo?'') sugiere la respuesta dentro de la propia pregunta y debe ser objetada. \\
\addlinespace[0.5em]

\textbf{¿Qué ocurre si dos abogados no se ponen de acuerdo sobre si una pregunta es indicativa?} &
Se retira al testigo, se discute la reformulación entre las partes y, de no haber acuerdo, el tribunal decide: reformula la pregunta él mismo o dispone que esa pregunta no se formule. \\
\addlinespace[0.5em]

\textbf{¿Son apelables las resoluciones sobre producción o denegación de prueba?} &
No, son irrecurribles e inapelables. La única vía disponible es el recurso de reposición, que debe interponerse y fundarse en el mismo acto de la audiencia, previo traslado a la contraria. \\
\addlinespace[0.5em]

\textbf{¿Cuáles son los distintos tipos de domicilio relevantes en un proceso judicial?} &
Domicilio real (donde vive la persona o está constituida la sociedad), denunciado (el que indica la contraria al demandar), constituido (fijado a los fines del proceso, dentro del radio del juzgado), electrónico (en el sistema informático judicial), fiscal (registrado ante el organismo recaudador) y legal (el que la ley asigna por función a determinadas personas). \\
\addlinespace[0.5em]

\textbf{¿Qué hacer si un cliente descubre un embargo originado en una notificación dirigida a un domicilio incorrecto?} &
Verificar primero el domicilio al que se dirigió la cédula. Si no coincide con el domicilio real (por ejemplo, el que consta en el contrato social ante la IGJ), corresponde plantear la nulidad de esa notificación y, como consecuencia, de la rebeldía y la sentencia dictada a partir de ella, solicitando el levantamiento del embargo. \\
\addlinespace[0.5em]

\textbf{¿El embargo en sí mismo es nulo en ese caso, o corresponde otro pedido distinto?} &
El embargo no es nulo en sí mismo, porque se trabó válidamente conforme a la sentencia vigente en ese momento. Lo que corresponde pedir, tras la nulidad de los actos anteriores, es su levantamiento. \\
\addlinespace[0.5em]

\textbf{¿Qué habilidad profesional destaca el docente como la más relevante hacia el futuro?} &
No la acumulación de conocimiento memorizado ---fácilmente accesible hoy con solo ``apretar un botón''--- sino la capacidad de manejar esa información con criterio, para no incurrir en el riesgo de creer resuelto un problema por el solo hecho de haber consultado un dato. \\

\midrule
\multicolumn{2}{p{0.94\textwidth}}{
\textbf{Resumen:} La clase reconstruye, de punta a punta, el camino de un hecho desde que se lo quiere probar hasta que una decisión judicial errónea sobre ese camino puede revertirse. Comienza revisando una audiencia testimonial real para mostrar que la teoría y la práctica se tocan: quién puede declarar, con qué requisitos y qué carga tiene la parte que lo ofrece. Avanza luego hacia la técnica concreta del interrogatorio ---preguntas abiertas frente a indicativas--- y hacia los recursos disponibles cuando una decisión sobre la prueba no conforma a una de las partes, con la irrecurribilidad como regla y la reposición en el acto como única vía. Cierra con un caso real que enlaza todo lo anterior desde otro ángulo: la validez de una notificación ---determinada, en última instancia, por el domicilio correcto--- es el punto de apoyo de toda la cadena procesal posterior; si ese punto de apoyo es defectuoso, caen la rebeldía, la sentencia y, en consecuencia, corresponde solicitar el levantamiento del embargo trabado a partir de ella. El hilo que conecta los cinco bloques es siempre el mismo: la eficacia de un litigante no depende solo de conocer la norma, sino de aplicarla con método y en el momento oportuno ---notificando a tiempo, objetando en el acto, verificando el domicilio antes de dar cualquier situación procesal por firme---.
} \\
\bottomrule
\end{longtable}

\end{document}
